\clearpage\pagestyle{matSimple}\phantomsection\label{reperage}\pdfbookmark[0]{Repérage du programme}{bm-reperage}
\matPageTitre{Repérage du programme}{Moyenne Section · objectifs « à partir de 4 ans » du programme de maternelle}
{\fontsize{9.5}{13.5}\selectfont Chaque objectif d’apprentissage du programme pour cette tranche d’âge est relié aux activités qui le travaillent. Les numéros renvoient aux pages du cahier.\par}\vspace{5mm}
{\fontsize{9.5}{13}\selectfont\begin{tabularx}{\linewidth}{@{}Xr@{}}\toprule
{\otSemi Objectif du programme} & {\otSemi Activités}\\\midrule
Comprendre qu’une quantité ne dépend ni de la nature des objets ni de leur organisation (jusqu’à six) & \hyperref[act-ms-1]{1}, \hyperref[act-ms-3]{3}\\[1.6mm]
Ajouter un : la quantité suivante, le nombre suivant & \hyperref[act-ms-10]{10}\\[1.6mm]
Parcourir une collection en passant une et une seule fois par chaque élément & \hyperref[act-ms-12]{12}\\[1.6mm]
Dénombrer une collection (jusqu’à six) & \hyperref[act-ms-4]{4}, \hyperref[act-ms-8]{8}\\[1.6mm]
Constituer une collection d’un cardinal donné (jusqu’à six) & \hyperref[act-ms-9]{9}\\[1.6mm]
Comparer des quantités : globalement, terme à terme, en dénombrant & \hyperref[act-ms-2]{2}, \hyperref[act-ms-5]{5}\\[1.6mm]
Composer et décomposer les nombres jusqu’à six & \hyperref[act-ms-13]{13}, \hyperref[act-ms-14]{14}\\[1.6mm]
Associer une quantité, le nom d’un nombre et son écriture chiffrée & \hyperref[act-ms-6]{6}, \hyperref[act-ms-7]{7}, \hyperref[act-ms-11]{11}\\[1.6mm]
Écrire en chiffres les nombres de un à six & \hyperref[act-ms-15]{15}, \hyperref[act-ms-16]{16}, \hyperref[act-ms-17]{17}\\[1.6mm]
Connaître la comptine numérique de un à douze & \hyperref[act-ms-18]{18}\\[1.6mm]
Comprendre la notion de rang (jusqu’au sixième) & \hyperref[act-ms-19]{19}, \hyperref[act-ms-20]{20}, \hyperref[act-ms-21]{21}\\[1.6mm]
Déterminer l’effet d’un déplacement : avancer ou reculer d’une ou deux cases & \hyperref[act-ms-24]{24}, \hyperref[act-ms-25]{25}\\[1.6mm]
Se familiariser avec le début de la bande numérique (jusqu’à six) & \hyperref[act-ms-22]{22}, \hyperref[act-ms-23]{23}\\[1.6mm]
Problèmes : rechercher le tout ou une partie & \hyperref[act-ms-26]{26}, \hyperref[act-ms-27]{27}, \hyperref[act-ms-28]{28}, \hyperref[act-ms-29]{29}\\[1.6mm]
Problèmes : groupements et partage équitable & \hyperref[act-ms-30]{30}, \hyperref[act-ms-31]{31}, \hyperref[act-ms-32]{32}\\[1.6mm]
Solides : produire et retrouver des empreintes & \hyperref[act-ms-35]{35}, \hyperref[act-ms-36]{36}\\[1.6mm]
Formes planes : reconnaître et classer triangles, carrés, disques & \hyperref[act-ms-33]{33}, \hyperref[act-ms-34]{34}\\[1.6mm]
Reproduire des assemblages de formes (au plus cinq) & \hyperref[act-ms-37]{37}, \hyperref[act-ms-38]{38}\\[1.6mm]
Longueurs : comparer, classer en trois groupes, ordonner & \hyperref[act-ms-39]{39}, \hyperref[act-ms-40]{40}, \hyperref[act-ms-41]{41}\\[1.6mm]
Masses : comparer deux objets, utiliser une balance de Roberval & \hyperref[act-ms-42]{42}, \hyperref[act-ms-43]{43}\\[1.6mm]
Motifs répétitifs : reconnaître, trouver l’intrus, prolonger, mémoriser & \hyperref[act-ms-44]{44}, \hyperref[act-ms-45]{45}, \hyperref[act-ms-46]{46}, \hyperref[act-ms-47]{47}, \hyperref[act-ms-48]{48}\\[1.6mm]
Se repérer dans l’espace : situer des objets, suivre et raconter un déplacement & \hyperref[act-ms-49]{49}, \hyperref[act-ms-50]{50}, \hyperref[act-ms-51]{51}, \hyperref[act-ms-52]{52}, \hyperref[act-ms-53]{53}\\[1.6mm]
\bottomrule\end{tabularx}\par}
